% Options for packages loaded elsewhere
\PassOptionsToPackage{unicode}{hyperref}
\PassOptionsToPackage{hyphens}{url}
\PassOptionsToPackage{dvipsnames,svgnames,x11names}{xcolor}
%
\documentclass[
  11pt,
]{article}
\usepackage{amsmath,amssymb}
\usepackage{iftex}
\ifPDFTeX
  \usepackage[T1]{fontenc}
  \usepackage[utf8]{inputenc}
  \usepackage{textcomp} % provide euro and other symbols
\else % if luatex or xetex
  \usepackage{unicode-math} % this also loads fontspec
  \defaultfontfeatures{Scale=MatchLowercase}
  \defaultfontfeatures[\rmfamily]{Ligatures=TeX,Scale=1}
\fi
\usepackage[]{charter}
\ifPDFTeX\else
  % xetex/luatex font selection
\fi
% Use upquote if available, for straight quotes in verbatim environments
\IfFileExists{upquote.sty}{\usepackage{upquote}}{}
\IfFileExists{microtype.sty}{% use microtype if available
  \usepackage[]{microtype}
  \UseMicrotypeSet[protrusion]{basicmath} % disable protrusion for tt fonts
}{}
\makeatletter
\@ifundefined{KOMAClassName}{% if non-KOMA class
  \IfFileExists{parskip.sty}{%
    \usepackage{parskip}
  }{% else
    \setlength{\parindent}{0pt}
    \setlength{\parskip}{6pt plus 2pt minus 1pt}}
}{% if KOMA class
  \KOMAoptions{parskip=half}}
\makeatother
\usepackage{xcolor}
\usepackage[margin=2.5cm]{geometry}
\setlength{\emergencystretch}{3em} % prevent overfull lines
\providecommand{\tightlist}{%
  \setlength{\itemsep}{0pt}\setlength{\parskip}{0pt}}
\setcounter{secnumdepth}{-\maxdimen} % remove section numbering
% ==========================================================================
%  IA na Sociedade — cabeçalho LaTeX (pdflatex)
%  Programa CiberExt 26-29 · FEELT38103
%  Adaptado de "AI in Society" (Universidade de Helsinque / Una Europa)
%  Licença: CC BY-NC-SA 4.0
% --------------------------------------------------------------------------
%  Espelha a identidade do ia-style.css: papel quente, tinta ardósia e a
%  faixa tricolor dos três eixos do curso (conceitos / aplicações / riscos).
% ==========================================================================

% ---------- Fontes ----------
% A serifada vem do build.sh, que passa  -V fontfamily=charter  ao pandoc.
% Aqui só acrescentamos a sem-serifa; ambas com fallback, porque nem toda
% instalação de LaTeX traz os mesmos pacotes.
\IfFileExists{helvet.sty}{\usepackage[scaled=0.90]{helvet}}{}
\usepackage{textcomp}
\IfFileExists{microtype.sty}{\usepackage{microtype}}{}

% ---------- Pacotes de layout ----------
\usepackage{xcolor}
\usepackage{tcolorbox}
\tcbuselibrary{skins,breakable}
\usepackage{titlesec}
\usepackage{fancyhdr}
\usepackage{enumitem}
\usepackage{amssymb}
\usepackage{booktabs}
\usepackage{array}
\usepackage{colortbl}
\usepackage{ragged2e}
\usepackage{fvextra}
\usepackage{newunicodechar}

% ---------- Caracteres Unicode que o pdflatex não conhece ----------
\newunicodechar{✓}{{\color{iaAplic}\checkmark}}
\newunicodechar{✅}{{\color{iaAplic}\checkmark}}
\newunicodechar{→}{\ensuremath{\rightarrow}}
\newunicodechar{←}{\ensuremath{\leftarrow}}
\newunicodechar{↔}{\ensuremath{\leftrightarrow}}
\newunicodechar{–}{--}
\newunicodechar{—}{---}
\newunicodechar{€}{\texteuro}
\newunicodechar{£}{\textsterling}
\newunicodechar{×}{\ensuremath{\times}}
\newunicodechar{≈}{\ensuremath{\approx}}
\newunicodechar{≥}{\ensuremath{\geq}}
\newunicodechar{≤}{\ensuremath{\leq}}
\newunicodechar{•}{\textbullet}
\newunicodechar{″}{''}
\newunicodechar{′}{'}

% ==========================================================================
%  PALETA — os mesmos valores do ia-style.css
% ==========================================================================
\definecolor{iaConceitos}{HTML}{24505C}   % petróleo — títulos, eixo 1
\definecolor{iaAplic}{HTML}{16746C}       % verdete  — subtítulos, eixo 2
\definecolor{iaRiscos}{HTML}{B45F14}      % açafrão  — destaques, eixo 3
\definecolor{iaAmeixa}{HTML}{77345F}      % links internos, termos técnicos
\definecolor{iaTinta}{HTML}{1F2933}       % corpo do texto
\definecolor{iaTintaFraca}{HTML}{5C6670}
\definecolor{iaRegra}{HTML}{E4DBCD}
\definecolor{iaBloco}{HTML}{F3EEE4}
\definecolor{iaNota}{HTML}{EAF1F0}
\definecolor{iaCaso}{HTML}{FCF3E6}
\definecolor{iaFilos}{HTML}{F7EFF4}

\color{iaTinta}

% ==========================================================================
%  ASSINATURA — faixa tricolor dos três eixos
% ==========================================================================
\newcommand{\faixaEixos}[1][\linewidth]{%
  \noindent
  \textcolor{iaConceitos}{\rule{0.3333#1}{4pt}}%
  \textcolor{iaAplic}{\rule{0.3333#1}{4pt}}%
  \textcolor{iaRiscos}{\rule{0.3333#1}{4pt}}%
  \par}

% ---------- Página de título ----------
\makeatletter
\newcommand{\subtitle}[1]{\gdef\@subtitle{#1}}
\providecommand{\@subtitle}{}
\def\@maketitle{%
  \newpage\null\vskip 1.5em%
  \faixaEixos
  \vskip 2.2em
  {\raggedright
    {\sffamily\bfseries\huge\color{iaConceitos}\@title\par}%
    \vskip 0.7em%
    {\sffamily\large\color{iaTintaFraca}\@subtitle\par}%
    \vskip 2em%
    {\color{iaRegra}\rule{\linewidth}{1pt}}%
    \vskip 1em%
    {\sffamily\small\color{iaTinta}\@author\par}%
    \vskip 0.5em%
    {\sffamily\footnotesize\color{iaTintaFraca}\@date\par}%
  }\par\vskip 2.5em}
\makeatother

% ==========================================================================
%  TÍTULOS DE SEÇÃO
% ==========================================================================
\titleformat{\section}[block]
  {\sffamily\Large\bfseries\color{iaConceitos}}
  {\thesection}{0.7em}{}
  [{\vspace{2pt}\color{iaRegra}\titlerule[1.2pt]}]
\titleformat{\subsection}
  {\sffamily\large\bfseries\color{iaAplic}}
  {\thesubsection}{0.6em}{}
\titleformat{\subsubsection}
  {\sffamily\normalsize\bfseries\color{iaRiscos}}
  {\thesubsubsection}{0.5em}{}
\titleformat{\paragraph}[runin]
  {\sffamily\small\bfseries\color{iaRiscos}}
  {}{0em}{}

\titlespacing*{\section}{0pt}{2.6ex plus 1ex minus .2ex}{1.3ex plus .2ex}
\titlespacing*{\subsection}{0pt}{2.2ex plus 1ex minus .2ex}{1ex plus .2ex}

% ==========================================================================
%  CABEÇALHO E RODAPÉ
% ==========================================================================
\setlength{\headheight}{24pt}
\renewcommand{\sectionmark}[1]{\markboth{#1}{}}
\pagestyle{fancy}
\fancyhf{}
\renewcommand{\headrulewidth}{0.6pt}
\renewcommand{\footrulewidth}{0pt}
\fancyhead[L]{\sffamily\footnotesize\color{iaAplic}Ética da IA}
\fancyhead[R]{\sffamily\footnotesize\color{iaAplic}\nouppercase{\leftmark}}
\fancyfoot[C]{\sffamily\footnotesize\color{iaConceitos}\thepage}
\renewcommand{\headrule}{%
  \hbox to\headwidth{\color{iaRegra}\leaders\hrule height \headrulewidth\hfill}}

% ==========================================================================
%  TEXTO CORRIDO
% ==========================================================================
\linespread{1.08}
\setlength{\parskip}{0.55em}
\setlength{\parindent}{0pt}
\emergencystretch=3em
\hbadness=10000

% ---------- Termos técnicos / código inline ----------
\let\oldtexttt\texttt
\renewcommand{\texttt}[1]{{\color{iaAmeixa}\small\ttfamily #1}}

% ---------- Blocos verbatim ----------
\fvset{breaklines=true,breakanywhere=true,fontsize=\small}

% ---------- Blocos de código (ambiente Shaded do pandoc) ----------
\renewenvironment{Shaded}
  {\begin{tcolorbox}[
      breakable, enhanced jigsaw, rounded corners,
      arc=3pt, boxrule=0pt, frame hidden,
      colback=iaBloco,
      borderline west={3pt}{0pt}{iaAmeixa},
      left=10pt, right=8pt, top=6pt, bottom=6pt,
      before skip=9pt, after skip=9pt]}
  {\end{tcolorbox}}

% ---------- Citações: caixa de caso ----------
\renewenvironment{quote}
  {\begin{tcolorbox}[
      breakable, enhanced jigsaw, rounded corners,
      arc=3pt, boxrule=0pt, frame hidden,
      colback=iaCaso,
      borderline west={3pt}{0pt}{iaRiscos},
      left=12pt, right=10pt, top=9pt, bottom=9pt,
      before skip=9pt, after skip=9pt]}
  {\end{tcolorbox}}

% ==========================================================================
%  BLOCOS DE APOIO
%  No markdown, escreva:   ::: nota   ...texto...   :::
%  O pandoc converte fenced_divs para os ambientes abaixo.
% ==========================================================================
\newtcolorbox{iabox}[3]{
  breakable, enhanced jigsaw, rounded corners,
  arc=3pt, boxrule=0pt, frame hidden,
  colback=#2,
  borderline west={3pt}{0pt}{#1},
  left=12pt, right=10pt, top=9pt, bottom=9pt,
  before skip=11pt, after skip=11pt,
  before upper={\setlength{\parskip}{0.5em}\setlength{\parindent}{0pt}},
  title={\sffamily\bfseries\footnotesize\textcolor{#1}{\MakeUppercase{#3}}},
  attach title to upper={\par\vspace{5pt}},
  fonttitle=\normalfont}

% Cada ambiente aceita um titulo opcional entre colchetes, vindo do
% atributo data-titulo do markdown:   \begin{filosofia}[Titulo] ... \end{filosofia}
% Os quatro primeiros espelham os icones do material original.
\newenvironment{filosofia}[1][Conceito filosofico]{\begin{iabox}{iaAmeixa}{iaFilos}{#1}}{\end{iabox}}
\newenvironment{tecnica}[1][Nota tecnica]{\begin{iabox}{iaConceitos}{iaNota}{#1}}{\end{iabox}}
\newenvironment{contexto}[1][Contexto]{\begin{iabox}{iaAplic}{iaBloco}{#1}}{\end{iabox}}
\newenvironment{definicao}[1][Definicao]{\begin{iabox}{iaRiscos}{iaCaso}{#1}}{\end{iabox}}

\newenvironment{nota}[1][Nota]{\begin{iabox}{iaConceitos}{iaNota}{#1}}{\end{iabox}}
\newenvironment{caso}[1][Caso]{\begin{iabox}{iaRiscos}{iaCaso}{#1}}{\end{iabox}}
\newenvironment{reflexao}[1][Para pensar]{\begin{iabox}{iaAplic}{iaBloco}{#1}}{\end{iabox}}
\newenvironment{licenca}
  {\par\vspace{2em}\faixaEixos\par\vspace{0.8em}
   \begingroup\sffamily\scriptsize\color{iaTintaFraca}}
  {\endgroup\par}

% ==========================================================================
%  LISTAS
% ==========================================================================
\setlist[itemize]{leftmargin=1.4em, itemsep=2pt, topsep=4pt, parsep=0pt}
\setlist[enumerate]{leftmargin=1.7em, itemsep=2pt, topsep=4pt, parsep=0pt}
\renewcommand{\labelitemi}{\textcolor{iaRiscos}{\textbullet}}
\renewcommand{\labelitemii}{\textcolor{iaAplic}{--}}

% ==========================================================================
%  TABELAS
% ==========================================================================
\renewcommand{\arraystretch}{1.3}
\arrayrulecolor{iaRegra}

% ==========================================================================
%  RÓTULOS EM PORTUGUÊS
%  Definidos aqui (e não via -V na linha de comando) porque acentos passados
%  como argumento dependem do locale do terminal e podem se corromper.
% ==========================================================================
\renewcommand{\contentsname}{Sumário}
\renewcommand{\refname}{Referências}
\renewcommand{\figurename}{Figura}
\renewcommand{\tablename}{Tabela}
\ifLuaTeX
  \usepackage{selnolig}  % disable illegal ligatures
\fi
\IfFileExists{bookmark.sty}{\usepackage{bookmark}}{\usepackage{hyperref}}
\IfFileExists{xurl.sty}{\usepackage{xurl}}{} % add URL line breaks if available
\urlstyle{same}
\hypersetup{
  pdftitle={Ética da Inteligência Artificial},
  colorlinks=true,
  linkcolor={iaAmeixa},
  filecolor={Maroon},
  citecolor={Blue},
  urlcolor={iaRiscos},
  pdfcreator={LaTeX via pandoc}}

\title{Ética da Inteligência Artificial}
\usepackage{etoolbox}
\makeatletter
\providecommand{\subtitle}[1]{% add subtitle to \maketitle
  \apptocmd{\@title}{\par {\large #1 \par}}{}{}
}
\makeatother
\subtitle{Capítulo 7 --- Exercícios}
\author{Tradução e adaptação para o português: José Lucas Lira Bizil,
Fernando Mazzeto Lisboa Lima e Matheus da Silva Fernandes\\
Programa CiberExt 26-29 · FEELT38103 · Universidade Federal de
Uberlândia\\
Original: \emph{Ethics of AI}, Universidade de Helsinque}
\date{2026}

\begin{document}
\maketitle

{
\hypersetup{linkcolor=iaConceitos}
\setcounter{tocdepth}{2}
\tableofcontents
}
\hypertarget{capuxedtulo-7-exercuxedcios}{%
\section{Capítulo 7 --- Exercícios}\label{capuxedtulo-7-exercuxedcios}}

\begin{nota}[Sobre estes exercícios]

Os três exercícios deste capítulo seguem o mesmo formato: leia a
citação, diga se concorda com suas afirmações e explique sua resposta.
Juntos, apresentam três posições distintas sobre a mesma questão --- se
a ética da IA tem ou não efeito prático.

\textbf{As citações foram traduzidas para o português.} Os textos
originais estão disponíveis nos links indicados, em inglês.

\end{nota}

\hypertarget{exercuxedcio-12a-a-uxe9tica-da-ia-e-vocuxea-dissertativa}{%
\subsection{\texorpdfstring{Exercício 12a --- A ética da IA e você
\texttt{{[}dissertativa{]}}}{Exercício 12a --- A ética da IA e você {[}dissertativa{]}}}\label{exercuxedcio-12a-a-uxe9tica-da-ia-e-vocuxea-dissertativa}}

\textbf{Pontuação:} 2 pontos · \textbf{Extensão:} 75 a 1500 palavras ·
\textbf{Tentativas:} ilimitadas \textbf{Componente Open edX:}
\texttt{openassessment} (revisão por pares)

Leia a citação a seguir. Você concorda com suas afirmações? Responda
brevemente e explique sua resposta.

\begin{quote}
Em 2019 falou-se mais de ética da IA do que nunca. Dezenas de
organizações produziram diretrizes de ética da IA; empresas correram
para constituir equipes de IA responsável e exibi-las diante da
imprensa. Já é difícil participar de uma conferência sobre IA sem que
parte da programação seja dedicada a alguma mensagem relacionada à
ética: como proteger a privacidade das pessoas quando a IA precisa de
tantos dados? Como empoderar comunidades marginalizadas em vez de
explorá-las? Como continuar a confiar na mídia diante da desinformação
criada e distribuída por algoritmos?

Mas conversa é só conversa --- não basta. Apesar de todo o discurso
dedicado a essas questões, as diretrizes de ética da IA de muitas
organizações permanecem vagas e difíceis de implementar. Poucas empresas
conseguem mostrar mudanças concretas na maneira como produtos e serviços
de IA são avaliados e aprovados. Estamos caindo na armadilha do
\emph{ethics-washing}, em que a ação genuína é substituída por promessas
superficiais\ldots{}

Karen Hao, ``In 2020, let's stop AI ethics-washing and actually do
something'', \emph{MIT Technology Review}.
\url{https://www.technologyreview.com/2019/12/27/57/ai-ethics-washing-time-to-act/}
\end{quote}

\begin{center}\rule{0.5\linewidth}{0.5pt}\end{center}

\hypertarget{exercuxedcio-12b-a-uxe9tica-da-ia-e-vocuxea-dissertativa}{%
\subsection{\texorpdfstring{Exercício 12b --- A ética da IA e você
\texttt{{[}dissertativa{]}}}{Exercício 12b --- A ética da IA e você {[}dissertativa{]}}}\label{exercuxedcio-12b-a-uxe9tica-da-ia-e-vocuxea-dissertativa}}

\textbf{Pontuação:} 2 pontos · \textbf{Extensão:} 75 a 1500 palavras ·
\textbf{Tentativas:} ilimitadas \textbf{Componente Open edX:}
\texttt{openassessment} (revisão por pares)

Leia a citação a seguir. Você concorda com suas afirmações? Responda
brevemente e explique sua resposta.

\begin{quote}
Lidar com o papel da filosofia e da ética exige superar ambas as
tendências e enxergar a ética como um modo de investigação que facilita
a avaliação de estratégias concorrentes de política tecnológica\ldots{}
Longe de impor um esquema autorregulatório ou determinada estrutura de
governança, a filosofia moral, na verdade, facilita o questionamento e a
reconsideração de qualquer prática dada, situando-a numa teia complexa
de instituições jurídicas, políticas e econômicas.

A filosofia moral pode, de fato, lançar nova luz sobre práticas humanas
ao acrescentar a perspectiva necessária, explicando a relação entre a
tecnologia e outros objetivos dignos, situando a tecnologia no humano,
no social, no político. Tornou-se urgente começar a considerar a ética
da tecnologia também de dentro, e não apenas de fora da ética.

Elena Bietti, ``From Ethics Washing to Ethics Bashing: A View on Tech
Ethics from Within Moral Philosophy''.
\url{https://papers.ssrn.com/sol3/papers.cfm?abstract_id=3513182}
\end{quote}

\begin{center}\rule{0.5\linewidth}{0.5pt}\end{center}

\hypertarget{exercuxedcio-12c-a-uxe9tica-da-ia-e-vocuxea-dissertativa}{%
\subsection{\texorpdfstring{Exercício 12c --- A ética da IA e você
\texttt{{[}dissertativa{]}}}{Exercício 12c --- A ética da IA e você {[}dissertativa{]}}}\label{exercuxedcio-12c-a-uxe9tica-da-ia-e-vocuxea-dissertativa}}

\textbf{Pontuação:} 2 pontos · \textbf{Extensão:} 75 a 1500 palavras ·
\textbf{Tentativas:} ilimitadas \textbf{Componente Open edX:}
\texttt{openassessment} (revisão por pares)

Leia a citação a seguir. Você concorda com suas afirmações? Responda
brevemente e explique sua resposta.

\begin{quote}
A ética tem dentes poderosos. Infelizmente, mal os estamos usando na
ética da IA --- não é de admirar, então, que a ética da IA seja chamada
de desdentada. Este artigo reflete sobre os diversos mecanismos éticos
que emergiram nos últimos anos em resposta à implantação e ao uso
massivos da IA na sociedade e aos riscos associados. Esses mecanismos
incluem listas de princípios, códigos de ética, recomendações e
diretrizes. Ainda assim, como muitos mostraram, esses desenvolvimentos
éticos, embora promissores, também são problemáticos: sua eficácia ainda
está por ser demonstrada, e são particularmente suscetíveis à
manipulação\ldots{}

Anaïs Rességuier e Rowena Rodrigues, ``AI ethics should not remain
toothless! A call to bring back the teeth of ethics''.
\url{https://journals.sagepub.com/doi/full/10.1177/2053951720942541}
\end{quote}

\begin{center}\rule{0.5\linewidth}{0.5pt}\end{center}

\hypertarget{avaliauxe7uxe3o-do-capuxedtulo-e-do-curso-pesquisa}{%
\subsection{\texorpdfstring{Avaliação do capítulo e do curso
\texttt{{[}pesquisa{]}}}{Avaliação do capítulo e do curso {[}pesquisa{]}}}\label{avaliauxe7uxe3o-do-capuxedtulo-e-do-curso-pesquisa}}

\textbf{Pontuação:} 1 ponto · \textbf{Tentativas:} ilimitadas

Desta vez a avaliação é também sobre o curso como um todo. Conte-nos o
que foi bom e o que poderia ser melhorado no futuro.

\begin{itemize}
\tightlist
\item
  \textbf{Utilidade do curso como um todo} --- nota de 1 a 5
\item
  \textbf{Clareza do curso como um todo} --- nota de 1 a 5
\item
  \textbf{Dificuldade do curso como um todo} --- nota de 1 a 5
\item
  \textbf{Comentários finais sobre o capítulo e o curso} --- até 1200
  palavras
\end{itemize}

\begin{center}\rule{0.5\linewidth}{0.5pt}\end{center}

\begin{licenca}

\textbf{Material original:} \emph{Ethics of AI}, Universidade de
Helsinque. Professores responsáveis: Anna-Mari Rusanen e Jukka K.
Nurminen. \url{https://ethics-of-ai.mooc.fi/}. \textbf{Esta versão:}
tradução por José Lucas Lira Bizil, Fernando Mazzeto Lisboa Lima e
Matheus da Silva Fernandes, Programa CiberExt 26-29 (FEELT38103), UFU,
2026. Obra derivada. \textbf{Licença:} CC BY-NC-SA 4.0.

\end{licenca}

\end{document}
